\documentclass[11pt]{article}

\usepackage{amsmath}
\usepackage{amssymb}
\usepackage[margin=1in]{geometry}

% Supporting-information Methods section describing the saturation dose--response
% estimator. Standalone-compilable (pdflatex); the section body can be lifted
% directly into a paper's supplement.

\title{Estimation of the temperature dose--response to aerosol release}
\author{}
\date{}

\begin{document}
\maketitle

\section{Experimental design and notation}

Near-surface air temperature ($T2$, in K) was obtained from an ensemble of
mesoscale simulations. The design crosses two meteorological \emph{episodes},
two spatial evaluation \emph{areas} (an urban ``city'' domain and a larger
``region'' domain), four \emph{release cases} (a no-release control and three
active cases), and $N$ stochastic \emph{ensemble members}, with output saved at
hourly resolution over $H$ hours ($H=121$, about five days). In the analyzed
data $N=3$.

Each release case is characterized by a release rate $r$ (kt\,h$^{-1}$): the
control has $r=0$, and the three active cases have $r\in\mathcal{R}=\{1,10,100\}$.
Crucially, the ensemble members are \emph{matched} across release cases: member
$e$ uses the same perturbed initial conditions in the control and in every active
case, so member $e$ realizes one common weather trajectory that differs between
cases only through the release. We denote by $T2_{r,e}(h)$ the temperature for
release rate $r$, member $e$, at hour $h$.

The estimation below is carried out \emph{independently} for each of the four
(episode, area) combinations and, within each, largely hour by hour. To keep the
notation light we describe a single (episode, area); all quantities should be read
as conditional on it.

\section{Baseline and paired anomaly}

Because the members share meteorology, each temperature can be written
schematically as
\begin{equation}
  T2_{r,e}(h) = \mu_r(h) + m_e(h) + \varepsilon_{r,e}(h),
  \label{eq:decomp}
\end{equation}
where $\mu_r(h)$ is the dose-dependent mean of interest, $m_e(h)$ is a large
member-specific meteorological offset common to \emph{all} release cases of member
$e$ at hour $h$, and $\varepsilon_{r,e}(h)$ is residual noise. The offset $m_e(h)$
inflates the apparent scatter of the raw temperatures but is common to control and
active cases; it is removed exactly by differencing each active case against the
control of the \emph{same} member.

We therefore define the baseline temperature as the control ensemble mean,
\begin{equation}
  T2_0(h) = \frac{1}{N}\sum_{e=1}^{N} T2_{0,e}(h),
  \qquad
  \widehat{\mathrm{SE}}\big[T2_0(h)\big] = \frac{s_0(h)}{\sqrt{N}},
  \label{eq:baseline}
\end{equation}
where $s_0(h)$ is the sample standard deviation (with $N-1$ denominator) of the
control members at hour $h$, and the \emph{paired anomaly} for each active case as
\begin{equation}
  d_{r,e}(h) = T2_{r,e}(h) - T2_{0,e}(h),
  \qquad r\in\mathcal{R}.
  \label{eq:anomaly}
\end{equation}
Substituting \eqref{eq:decomp} shows that the common $m_e(h)$ cancels in
\eqref{eq:anomaly}, leaving the release signal plus a small differenced residual.
The baseline is taken directly from \eqref{eq:baseline} rather than fitted, so it
cannot trade against the dose--response curve.

\section{Saturation dose--response model}

For each hour we model the paired anomaly as a per-hour amplitude times a
saturating function of the release rate,
\begin{equation}
  d_{r,e}(h) \;=\; A(h)\, g(r;\rho) \;+\; \text{noise},
  \qquad
  g(r;\rho) = \frac{1 - e^{-r/\rho}}{\,1 - e^{-10/\rho}\,}.
  \label{eq:model}
\end{equation}
The shape function $g$ has two convenient properties: $g(0;\rho)=0$, so the model
passes through the origin (no anomaly without release), and $g(10;\rho)=1$, so the
amplitude $A(h)$ is exactly the temperature perturbation at $10$~kt\,h$^{-1}$
(reported as $\Delta T2_{10}$). The single shape parameter $\rho$ (the
\emph{release scale}, in kt\,h$^{-1}$) controls the curvature: as
$\rho\!\to\!\infty$ the response becomes linear in $r$, while as $\rho\!\to\!0$ it
saturates immediately; $\rho$ is the release rate at which the response reaches
$1-1/e\approx 63\%$ of its far-field value. The amplitude $A(h)$ is free at every
hour, whereas $\rho$ is held \emph{constant across all hours} within an
(episode, area).\footnote{An alternative power-law shape,
$g(r;\beta)=(r/10)^{\beta}$, was also fit within the identical framework; it shares
the properties $g(0)=0$, $g(10)=1$ and the same estimation and uncertainty
machinery, with $\beta$ in place of $\rho$.}

\section{Estimation by profiled least squares}

Conditional on $\rho$, model \eqref{eq:model} is linear in the per-hour amplitude,
so each $A(h)$ has a closed-form least-squares solution through the origin:
\begin{equation}
  \widehat{A}(h;\rho)
  = \frac{\sum_{r\in\mathcal{R}}\sum_{e=1}^{N} g(r;\rho)\, d_{r,e}(h)}
         {\sum_{r\in\mathcal{R}}\sum_{e=1}^{N} g(r;\rho)^2}.
  \label{eq:Ahat}
\end{equation}
The shape parameter is then estimated by profiling: $\widehat{\rho}$ minimizes the
total residual sum of squares obtained after concentrating out the amplitudes,
\begin{equation}
  \widehat{\rho}
  = \arg\min_{\rho}\;
    \sum_{h=1}^{H}\sum_{r\in\mathcal{R}}\sum_{e=1}^{N}
    \Big( d_{r,e}(h) - \widehat{A}(h;\rho)\, g(r;\rho) \Big)^2 .
  \label{eq:profile}
\end{equation}
This reduces the nonlinear fit to a one-dimensional search over the scalar $\rho$
wrapping the linear step \eqref{eq:Ahat}; the minimization is performed over
$\log\rho$ for scale invariance. The reported amplitudes are
$\widehat{A}(h) = \widehat{A}(h;\widehat{\rho})$.

\section{Uncertainty: the ensemble member as the unit of replication}
\label{sec:uncertainty}

Because the members are independent, matched realizations, each member supplies a
complete paired data set and hence one independent estimate of the amplitude. For
fixed $\rho$ define the per-member through-origin slope
\begin{equation}
  u_e(h;\rho)
  = \frac{\sum_{r\in\mathcal{R}} g(r;\rho)\, d_{r,e}(h)}
         {\sum_{r\in\mathcal{R}} g(r;\rho)^2}.
\end{equation}
The pooled estimate \eqref{eq:Ahat} equals the member mean of these slopes,
$\widehat{A}(h) = N^{-1}\sum_e u_e(h;\widehat{\rho})$, so its standard error
conditional on $\rho$ is the standard error of that mean,
\begin{equation}
  \widehat{\mathrm{SE}}\big[\widehat{A}(h)\,\big|\,\rho\big]
  = \frac{1}{\sqrt{N}}
    \left[\frac{1}{N-1}\sum_{e=1}^{N}\big(u_e(h;\rho)-\overline{u}(h)\big)^2\right]^{1/2},
  \qquad \overline{u}(h)=\frac{1}{N}\sum_{e} u_e(h;\rho).
  \label{eq:se_cond}
\end{equation}
Estimate \eqref{eq:se_cond} is exact and distribution-free: it makes no
homoscedasticity or additivity assumption and automatically incorporates the
correlation induced by the shared control, because the entire difference
\eqref{eq:anomaly} is contained within each $u_e$.

The uncertainty in the shape parameter $\widehat{\rho}$, and the amplitude
uncertainty that includes it, are obtained by a delete-one-member jackknife.
Removing member $e$ in its entirety---its control \emph{and} its active
cases---preserves the pairing; refitting \eqref{eq:profile}--\eqref{eq:Ahat} on the
remaining $N-1$ members yields $\widehat{\rho}_{(-e)}$ and
$\widehat{A}_{(-e)}(h)$. For a statistic $\theta$ (either $\widehat{\rho}$ or
$\widehat{A}(h)$),
\begin{equation}
  \widehat{\mathrm{SE}}_{\mathrm{jack}}[\theta]
  = \left[\frac{N-1}{N}\sum_{e=1}^{N}
     \big(\theta_{(-e)}-\overline{\theta}_{(\cdot)}\big)^2\right]^{1/2},
  \qquad
  \overline{\theta}_{(\cdot)}=\frac{1}{N}\sum_{e=1}^{N}\theta_{(-e)}.
  \label{eq:jack}
\end{equation}
The jackknife amplitude standard error includes the contribution of $\rho$
uncertainty; comparing it with the conditional standard error \eqref{eq:se_cond}
quantifies how much the shared shape parameter adds. All standard errors scale as
$N^{-1/2}$, so they contract as ensemble members are added.

\paragraph{Reading uncertainty bands.} Shaded intervals in the figures are
$\pm 1$ standard error. A nominal $95\%$ confidence interval multiplies the
standard error by a Student-$t$ factor with $N-1$ degrees of freedom; for $N=3$
this factor is $t_{2,\,0.975}\approx 4.30$, so a $95\%$ interval is roughly four
times wider than the plotted band. Both the multiplier (toward $1.96$) and the
standard errors themselves decrease as $N$ grows.

\section{Implementation}

Analyses used Python with NumPy, pandas, and SciPy. The one-dimensional profile
minimization \eqref{eq:profile} used \texttt{scipy.optimize.minimize\_scalar}
(bounded Brent method) over $\log\rho$ on $[\log 0.05,\ \log(2\times10^{4})]$. The
factor $1-e^{-x}$ was evaluated as $-\,\mathrm{expm1}(-x)$ to avoid loss of
precision at large $\rho$. The procedure was applied separately to each of the
four (episode, area) combinations, yielding one $\widehat{\rho}$ per combination,
per-hour amplitudes $\widehat{A}(h)=\Delta T2_{10}(h)$, and the baseline
$T2_0(h)$ from \eqref{eq:baseline}.

\section{Daily-maximum temperature reduction}
\label{sec:tmax}

The same estimator is applied to the effect of the release on the \emph{daily
high} temperature. Within each (episode, area) the hourly series is divided into
$D=5$ consecutive 24-hour blocks; the first hour is discarded so that hours
$2,\dots,121$ form five complete blocks. Let $\mathcal{H}_d$ index the hours of
block $d$. The daily maximum for release rate $r$, member $e$, day $d$ is
\begin{equation}
  M_{r,e,d} = \max_{h\in\mathcal{H}_d} T2_{r,e}(h),
\end{equation}
and the paired daily-maximum reduction is
\begin{equation}
  R_{r,e,d} = M_{0,e,d} - M_{r,e,d}.
  \label{eq:reduction}
\end{equation}
Equation \eqref{eq:reduction} is written control-minus-experiment, so that a
lowering of the daily high is a \emph{positive} reduction---the opposite sign
convention to the hourly anomaly \eqref{eq:anomaly}. The difference is taken
within a member, so the two maxima may fall at different hours and the shared
member meteorology still largely cancels.

The reduction is fit with the same saturation shape as \eqref{eq:model},
\begin{equation}
  R_{r,e,d} = A\, g(r;\rho),
  \qquad
  g(r;\rho) = \frac{1 - e^{-r/\rho}}{\,1 - e^{-10/\rho}\,},
  \label{eq:tmax_model}
\end{equation}
through the origin over $r\in\{1,10,100\}$. Here the amplitude $A$ (the reduction
at $10$~kt\,h$^{-1}$) and the release scale $\rho$ are \emph{scalars} for the
fitting group rather than hour-indexed. Two groupings are reported:
\begin{itemize}
  \item a \textbf{pooled} fit combining all days, giving one $(A,\rho)$ per
        (episode, area);
  \item a \textbf{per-day} fit, giving one $(A,\rho)$ per (episode, area, day).
\end{itemize}
Estimation and uncertainty are otherwise unchanged. For fixed $\rho$ the amplitude
is the closed-form through-origin slope $\widehat{A}=\sum g R / \sum g^2$
(Eq.~\eqref{eq:Ahat} with the per-hour index removed and $R$ in place of $d$);
$\rho$ is obtained by the same log-space profile minimization
\eqref{eq:profile}; the amplitude standard error conditional on $\rho$ follows
from the per-member through-origin slopes \eqref{eq:se_cond}; and $\rho$ together
with the $\rho$-inclusive amplitude standard error come from the
delete-one-member jackknife \eqref{eq:jack}. The baseline daily high itself,
$M_{0,e,d}$ averaged over members, is reported separately with its ensemble
standard error.

\paragraph{Identifiability of the per-day scale.} The pooled fit uses
$3\times3\times5 = 45$ paired reductions and yields a well-determined $\rho$; the
per-day fit uses only $3\times3 = 9$ and frequently cannot resolve the curvature,
with $\widehat{\rho}$ running toward the upper search bound on days whose
dose--response is effectively linear. The per-day \emph{amplitude} $A$ remains
well behaved and is the quantity to interpret day by day; the per-day \emph{scale}
should be read together with its (often large) jackknife interval, or $\rho$ held
common across days.

\section{Joint day--dose model for the daily maximum}

The pooled and per-day fits of Section~\ref{sec:tmax} sit at two extremes: the
pooled fit ignores any systematic evolution of the reduction over the episode,
while the per-day fit spends a separate amplitude \emph{and} scale on each day and
so resolves neither well. The joint model interpolates between them by letting the
reduction \emph{build up} over the episode with its own time scale, while sharing
one amplitude and one dose curvature across all days.

\subsection{Model}

The paired daily-maximum reduction $R_{r,e,d}$ of \eqref{eq:reduction} is modelled
as a product of a day factor and a dose factor,
\begin{equation}
  R_{r,e,d} \;=\; A_\infty\, f(d;\tau)\, g(r;\rho) \;+\; \text{noise},
  \label{eq:daydose}
\end{equation}
with the dose factor $g$ exactly as in \eqref{eq:tmax_model} and the day factor
\begin{equation}
  f(d;\tau) \;=\; 1 - e^{-(d-0.5)/\tau}.
  \label{eq:dayterm}
\end{equation}
Three parameters are estimated per (episode, area), giving four regressions in
total, each fit to all $3\times3\times5=45$ paired reductions:
\begin{itemize}
  \item $A_\infty$ (K), the \emph{day-asymptotic amplitude}: the reduction at
        $10$~kt\,h$^{-1}$ once the buildup has fully saturated ($d\to\infty$,
        $f\to1$). Because $g(10;\rho)=1$ by construction, $A_\infty$ is read
        directly in kelvin. It is reported as \texttt{Tmax\_scale10}.
  \item $\tau$ (days), the \emph{day scale}: the e-folding time of the buildup.
        $\tau\to0$ places the full reduction on day~1; $\tau\gg D$ makes the
        reduction grow nearly linearly across the episode.
  \item $\rho$ (kt\,h$^{-1}$), the release scale, with the same meaning as in
        \eqref{eq:model}.
\end{itemize}

The offset $-0.5$ in \eqref{eq:dayterm} places the origin of the buildup half a day
before the first daily maximum, so that $f$ vanishes there rather than at $d=1$;
day~1 therefore carries a partial response $f(1;\tau)=1-e^{-0.5/\tau}$ instead of
being forced to zero, which would be physically wrong for a release that is already
active during the first block. Days are integers $d\in\{1,\dots,5\}$ from the
blocking of Section~\ref{sec:tmax}.

Note that \eqref{eq:daydose} is \emph{separable}: the day factor multiplies the
dose factor, so the shape of the dose--response is assumed identical on every day
and only its overall size grows with $d$. Equivalently, the model asserts that the
release scale $\rho$ does not itself drift over the episode.

\subsection{Estimation}

For fixed $(\tau,\rho)$ the model \eqref{eq:daydose} is linear in $A_\infty$.
Writing the design value of each observation as
$x_{r,e,d}=f(d;\tau)\,g(r;\rho)$, the amplitude has the closed-form through-origin
solution
\begin{equation}
  \widehat{A}_\infty(\tau,\rho)
  = \frac{\sum_{r,e,d} x_{r,e,d}\, R_{r,e,d}}{\sum_{r,e,d} x_{r,e,d}^{2}},
  \label{eq:amp_joint}
\end{equation}
so only the two scales are searched. They minimize the residual sum of squares
after concentrating out the amplitude,
\begin{equation}
  (\widehat{\tau},\widehat{\rho})
  = \arg\min_{\tau,\rho}\;
    \sum_{r,e,d}\Big(R_{r,e,d} - \widehat{A}_\infty(\tau,\rho)\,f(d;\tau)\,g(r;\rho)\Big)^2 .
  \label{eq:profile_joint}
\end{equation}

The search is carried out over $(\log\tau,\log\rho)$ for scale invariance, by
Nelder--Mead simplex, from a $4\times4$ grid of starting points --- $\tau_0 \in
\{0.3, 1, 3, 30\}$~days crossed with $\rho_0 \in \{1, 10, 100, 1000\}$~kt\,h$^{-1}$
--- with the lowest attained objective retained. The multi-start guards against the
local minima that the ridge described below can induce. Because Nelder--Mead is
unconstrained, the objective clamps its argument into the box
\begin{equation}
  \log\tau \in [\log 0.02,\ \log 10^{3}], \qquad
  \log\rho \in [\log 0.05,\ \log(2\times10^{4})],
\end{equation}
and the returned optimum is clipped to the same box, so a reported optimum is
always attainable. Convergence used \texttt{xatol}~$=10^{-4}$,
\texttt{fatol}~$=10^{-10}$, and at most $2000$ iterations. As elsewhere, factors of
the form $1-e^{-x}$ were evaluated as $-\,\mathrm{expm1}(-x)$.

The coefficient of determination is $r^2 = 1 - \mathrm{SSR}/\mathrm{SST}$ with
$\mathrm{SST}$ taken about the mean of the observed reductions.

\subsection{Reported quantity: the day-5 reduction}

Because $A_\infty$ is an extrapolation to $d\to\infty$ --- beyond the five days
actually simulated --- the reduction on the final simulated day,
\begin{equation}
  R_{10,\,d=5} \;=\; \widehat{A}_\infty\, f(5;\widehat{\tau}),
  \label{eq:day5}
\end{equation}
is also reported (\texttt{Tmax\_red10\_day5}) and is the quantity to quote whenever
the day scale is poorly determined. It is an interpolation within the simulated
range and stays well determined even when $A_\infty$ and $\tau$ individually do
not.

\subsection{Uncertainty}

Uncertainty follows the same scheme as Section~\ref{sec:uncertainty}, with the
ensemble member again the unit of replication. Conditional on both scales, the
per-member through-origin slope is
\begin{equation}
  u_e = \frac{\sum_{r,d} x_{r,e,d}\, R_{r,e,d}}{\sum_{r,d} x_{r,e,d}^{2}},
\end{equation}
and the amplitude standard error conditional on $(\tau,\rho)$ is the standard error
of the mean of $\{u_e\}$, exactly as in \eqref{eq:se_cond} (reported as
\texttt{Tmax\_scale10\_se\_cond}).

Standard errors for $\widehat{\tau}$ and $\widehat{\rho}$, the $\rho$- and
$\tau$-inclusive standard error of $\widehat{A}_\infty$, and the standard error of
the day-5 reduction \eqref{eq:day5} all come from the delete-one-member jackknife
\eqref{eq:jack}. Removing member $e$ removes its control and all of its active
cases together, preserving the pairing, and the full three-parameter fit
\eqref{eq:amp_joint}--\eqref{eq:profile_joint} --- including the complete
multi-start search --- is repeated on the remaining $N-1$ members. For the day-5
reduction the jackknife replicate is formed as
$\widehat{A}_{\infty,(-e)}\, f(5;\widehat{\tau}_{(-e)})$, so that correlation
between the amplitude and the day scale is carried through rather than ignored.

\subsection{Identifiability of the day scale}

With only $D=5$ days the day scale is not always resolvable, and the failure is
structural rather than numerical. For $\tau \gg D$ the day factor linearizes,
\begin{equation}
  f(d;\tau) \;=\; 1-e^{-(d-0.5)/\tau} \;\approx\; \frac{d-0.5}{\tau},
\end{equation}
so that \eqref{eq:daydose} depends on the two parameters only through the ratio
$A_\infty/\tau$. Any $(A_\infty,\tau)$ sharing that ratio fits equally well: the
likelihood has a flat ridge, and $A_\infty$ and $\tau$ are individually
unidentified while their ratio --- and hence the fitted reduction on any simulated
day --- remains perfectly well constrained. The opposite extreme, $\tau$ at its
lower bound, means no buildup is resolved at all: the day factor is effectively
flat at unity and the model collapses to the pure dose saturation
\eqref{eq:tmax_model}.

Rather than let this pass silently, the fit records whether $\log\widehat{\tau}$
has landed on either end of its search box (within $10^{-3}$ in log units), as
\texttt{day\_scale\_at\_bound} $\in$ \{\texttt{lower}, \texttt{upper}\}. A pegged
scale also produces a degenerate jackknife standard error --- often exactly zero,
because every delete-one replicate is pinned at the same bound --- which must not
be read as precision. When \texttt{day\_scale\_at\_bound} is non-empty, $A_\infty$
and $\tau$ should not be interpreted separately, and the day-5 reduction
\eqref{eq:day5} should be quoted in place of $A_\infty$.

\end{document}
